\documentclass[11pt]{article}
\usepackage[a4paper,margin=2.6cm]{geometry}
\usepackage[T1]{fontenc}
\usepackage[utf8]{inputenc}
\usepackage{lmodern}
\usepackage{amsmath,amssymb,amsthm}
\usepackage{microtype}
\usepackage[hidelinks]{hyperref}
\usepackage{enumitem}
\usepackage{booktabs}
\setlist{itemsep=2pt,topsep=4pt}

\newtheorem*{result}{Result}
\newtheorem*{wip}{Work in progress}
\newcommand{\R}{\mathbb{R}}

\title{\vspace{-1.2cm}Progress on the maximum size of an almost-equidistant
set in $\R^5$:\\ $f(5)\le 17$, toward the exact value\\[4pt]
\large A computer-assisted campaign, with certified non-realizability at
$18$, $19$ and $20$ points}
\author{}
\date{August 5, 2026, third snapshot (level 17: 12615/12654; controls suite passed)}

\begin{document}
\maketitle
\vspace{-1.2cm}

\section*{Background}

A finite point set in $\R^d$ is \emph{almost equidistant} if among any three
of its points some two are at unit distance; $f(d)$ denotes the maximum size
of such a set. The known values are $f(1)=4$, $f(2)=7$, $f(3)=10$
\cite{Gyorey,BPSSV} and $f(4)=12$ (computer-assisted resolution of
Conjecture~1 of \cite{BPSSV}; see the accompanying note
\texttt{f4\_equals\_12.pdf} in this repository). For $d=5$ the published
state of the art is
\[
  16 \;\le\; f(5) \;\le\; 20,
\]
the lower bound by Larman--Rogers \cite{LR} (the odd half-cube
$\{v\in\{\pm1\}^5:\ v\ \text{has an odd number of}\ {+}1\text{s}\}/\sqrt8$,
whose non-unit-distance graph is the Clebsch graph), the upper bound by
Balko, P\'or, Scheucher, Swanepoel and Valtr \cite{BPSSV}. We found no
later improvement in the literature (2020--2026); the asymptotic bound is
$f(d)=O(d^{4/3})$ \cite{KMS}.

This note records the state of an ongoing computer-assisted campaign to
determine $f(5)$ exactly, run with the same certified machinery that proved
$f(4)=12$, generalized from $\R^4$ to $\R^5$. All code and data are in this
repository.

\begin{result}
$f(5) \le 17$. More precisely: there is no almost-equidistant set of $18$
points in $\R^5$ (and none of $19$ or $20$ points), by certified
elimination of all candidate unit-distance graphs at each of the three
levels. Combined with the Larman--Rogers set, $16 \le f(5) \le 17$.
\end{result}

\begin{wip}
One level remains: of the $12654$ candidate graphs on $17$ vertices,
$12615$ are already certified non-realizable at the time of this snapshot
($12518$ of them by circle-free searches taking six seconds in total);
the remaining $39$, all with a circle stage, exhibit narrow arcs of the
circle parameter on which full-width interval evaluation blows up ---
direct probes show every such arc dies within ${\sim}65$ search nodes
once split, so these are scheduling casualties, not geometric
obstructions, and are queued for a re-run with direct fine tiling
(\texttt{results\_n17.json} is updated continuously). The full
validation-control suite passes on the final kernel build (kill controls
$K_7$ and $K_1{+}K_{2,2,2,2,2}$ die in $1$--$2$ nodes; the sliced
cross-polytope survive control isolates the true mirror realizations in
exactly the $6/24$ slices containing them). Independent
numerical optimization over \emph{all} $12654$ candidates
(Levenberg--Marquardt, up to $300$ restarts per graph) found no graph
anywhere near realizable: the smallest residual is $0.409$, against
$10^{-30}$ for realizable controls. If every $17$-vertex candidate is
eliminated, $f(5)=16$; a certified survivor would instead pin $f(5)$
at $17$.
\end{wip}

\section*{Reduction}

Following \cite{BPSSV}, if $P\subset\R^5$ is almost equidistant then its
unit-distance graph $G$ satisfies: (i) $\alpha(G)\le 2$ (the complement is
triangle-free --- this is almost-equidistance itself); (ii) $G$ has no
$K_7$ (at most $d{+}1$ points of $\R^d$ are pairwise at unit distance);
(iii) $G$ has no $K_{3,3,3}$ subgraph (\cite{BPSSV}, Lemma~11: for odd $d$,
$K_{(d+1)/2}(3,\dots,3)$ is not a unit-distance graph in $\R^d$; three
$3$-point classes with all cross distances $1$ would need three pairwise
orthogonal $\ge2$-dimensional affine hulls, i.e.\ dimension $\ge 6$).
Graphs with (i)--(iii) are the \emph{abstract almost-equidistant graphs in
$\R^5$}.

Deleting edges of $G$ that are removable without violating (i) yields a
spanning \emph{minimal} abstract almost-equidistant subgraph $G'$
((ii),(iii) survive edge deletion). Any realization of $G$ restricts to
one of $G'$: $n$ distinct points with every $G'$-edge at distance exactly
$1$ (non-edges unconstrained). Hence: \emph{if no minimal abstract
almost-equidistant graph on $n$ vertices is realizable with $n$ distinct
points, then $f(5)<n$.} Minimal graphs are exactly the complements of
\emph{maximal triangle-free} graphs satisfying (ii) and (iii).

\section*{Candidate generation, independently verified}

Maximal triangle-free graphs were generated with \texttt{triangleramsey}
\cite{BGS,BBH} and filtered by our own checkers for (ii) and (iii)
(\texttt{filter\_mtf.c}); every input graph is re-verified to be maximal
triangle-free from scratch. The resulting counts reproduce
\cite{BPSSV} (Table~2, $d=5$ column) exactly:

\begin{center}
\begin{tabular}{lrrrrr}
\toprule
$n$ & 17 & 18 & 19 & 20 & 21\\
\midrule
maximal triangle-free graphs & 164\,796 & 1\,337\,848 & 13\,734\,745 &
178\,587\,364 & 2\,911\,304\,940\\
minimal abstract a.e.\ graphs & 12\,654 & 8\,825 & 340 & 8 & 0\\
\bottomrule
\end{tabular}
\end{center}

The $n=21$ row is an independent re-verification of $f(5)\le 20$
\cite{BPSSV}: all $2.9$ billion maximal triangle-free graphs on $21$
vertices were regenerated and none passes the filter.

As an independent check of the entire pipeline, a from-scratch enumerator
(\texttt{enumaeq5.c}, growing graphs vertex-by-vertex with incremental
tests of (i)--(iii), minimum-degree normalization, and isomorphism
rejection --- the method of Appendix~A of the $f(4)$ note) reproduces the
$d=5$ column of \cite{BPSSV} Table~3 at every order it reaches
($7$, $14$, $38$, $106$, $402$, $1817$, $11\,132$, $86\,053$, $803\,299$,
$7\,623\,096$ for $n=4,\dots,13$), and its minimal counts
($2,3,4,6,9,14,25,46,106,242$) match Table~2. At $n=13$ the $242$ minimal
graphs from the two independent routes agree \emph{as sets}, matched
one-to-one up to isomorphism (\texttt{compare\_sets.py}).

The lower bound is re-verified in integer arithmetic
(\texttt{verify\_lower\_bound\_16.py}): the $16$ half-cube points are
pairwise at squared distance $8$ or $16$ (unscaled), any three contain a
pair at $8$, and the non-unit graph is the $5$-regular triangle-free
Clebsch graph; scaling by $1/\sqrt8$ gives an almost-equidistant $16$-point
set, so $f(5)\ge 16$.

\section*{The certified decision engine in $\R^5$}

\texttt{ckernel5.c} is a direct port of the certified $\R^4$ engine
(\texttt{ckernel.c}; Appendix~B of the $f(4)$ note), with the
dimension-forced changes and nothing else:
\begin{itemize}
\item points are $5$-dimensional interval vectors (outward-rounded IEEE
      doubles; $\cos/\sin$ padded by $8$ ulp);
\item seeds are $K_6$ (canonical unit $5$-simplex, exact interval
      enclosures, rigid) or $K_5$ (mirror explored by branching);
      $12\,652$ of the $12\,654$ graphs on $17$ vertices contain $K_6$;
\item a vertex with $\ge5$ placed neighbours is placed via the
      $5$-sphere system ($4$ difference equations, $5$ unknowns,
      $1$-dimensional nullspace, quadratic --- at most two roots),
      contracted by a $5{\times}5$ Krawczyk operator on the
      best-conditioned $5$-subset of neighbours, all others acting as
      certified pruning constraints;
\item a vertex with exactly $4$ placed neighbours lies on a circle
      ($3{\times}3$ interval Gram system for the circumcentre,
      $2$-dimensional orthogonal complement basis by interval
      Gram--Schmidt), subdivided adaptively with the same runtime canary;
\item near-tangent (merged-root) cases are resolved by the same
      one-dimensional segment subdivision;
\item the two distinctness rules discard branches on which coincidence of
      non-adjacent vertices is forced: root separation now uses
      $5$-subsets of common placed neighbours; the bisector-hyperplane
      rule needs $6$ common placed neighbours with certifiably nonzero
      affine volume ($5{\times}5$ interval determinant).
\end{itemize}
A branch is discarded only when an interval certifiably excludes a
required value or when coincidence (fewer than $n$ distinct points) is
forced; eliminating all branches of one elimination order proves
non-realizability with $n$ distinct points. Elimination orders are
generated circle-late; for the dense graphs at $n\ge18$ almost all
searches are circle-free binary trees ($8813/8825$ at $n=18$; $339/340$ at
$n=19$; $7/8$ at $n=20$), and every one of the $21\,827$ candidate graphs
at levels $17$--$20$ admits a valid order (no two-parameter stage is ever
needed).

\emph{Controls.} Verified at the time of this snapshot: the engine kills
$K_7$ (no seven pairwise-unit points in $\R^5$) and does \emph{not} kill
the realizable controls $K_6$ (the unit $5$-simplex) and $K_{2,2,2,2,2}$
(the unit cross-polytope, recovered numerically with distinct points and
not killable by the engine within large node budgets). The remaining suite
--- the sliced survive controls for the cross-polytope and the $16$-vertex
half-cube graph, and the kill control $K_1+K_{2,2,2,2,2}$ (an apex joined
to the cross-polytope graph is non-realizable: the apex would have to be
the centre, at distance $1/\sqrt2\ne1$) --- is running at snapshot time;
its outcome will be recorded in \texttt{controls5\_run.log} and in a
revised note. (Survive controls run as sliced tilings and pass when slices
report floor-width survivors; kill controls require every slice to die.)

\section*{Certified results at $20$ and $19$ points}

For each graph, the driver (\texttt{reproduce5.py}) races the candidate
elimination orders concurrently; a graph is certified when one order's
tiling of the first circle parameter (a single unrestricted search for
circle-free orders) is entirely killed.

\begin{itemize}
\item \textbf{$n=20$:} all $8$ minimal abstract almost-equidistant graphs
      certified non-realizable with $20$ distinct points. Wall time
      ${\sim}5$ seconds total (Ryzen 9950X3D, LLVM clang build).
      Consequence: no $20$-point almost-equidistant set in $\R^5$, i.e.\
      $f(5)\le 19$.
\item \textbf{$n=19$:} all $340$ graphs certified non-realizable with
      $19$ distinct points (${\sim}170$ seconds wall under load).
      Consequence: $f(5)\le 18$.
\item \textbf{$n=18$:} all $8825$ graphs certified non-realizable with
      $18$ distinct points ($8813$ circle-free searches bulk-certified in
      ${\sim}10$ seconds of wall time on $20$ workers; the $12$ graphs
      with a circle stage certified by the racing driver in ${\sim}1$
      second each). Consequence: \textbf{$f(5)\le 17$}.
\end{itemize}

Per-graph verdicts and the winning decomposition indices are recorded in
\texttt{results\_n20.json}, \texttt{results\_n19.json} and
\texttt{results\_n18.json}. Numerical corroboration: Levenberg--Marquardt
finds best residuals $\ge 2.6$ at $n=19,20$ and $\ge 0.4$ over all of
$n=17$ (distinct-point solutions), i.e.\ nowhere near realizable.

\section*{Caveats}

The proof of $f(5)\le18$ is computer-assisted, with the same trust base as
the $f(4)=12$ note: IEEE-754 correctly rounded arithmetic, library
$\cos/\sin$ accurate to $8$ ulp (padded), the engine logic of Appendix~B
of that note with the $\R^5$ modifications above, and --- for the
\emph{completeness} of the candidate lists --- the \texttt{triangleramsey}
generator \cite{BGS}, cross-validated here against \cite{BPSSV} at
$n=17,\dots,21$ and against a fully independent enumeration up to $n=13$
(counts at every order, sets at $n=13$). The numerical sweeps are
heuristic corroboration only; no conclusion rests on them. A reproduction
run is one command per level (\texttt{python reproduce5.py --level 20},
etc.); the kernel compiles automatically and validation controls run
first.

\begin{thebibliography}{9}\small
\bibitem{BPSSV} M.~Balko, A.~P\'or, M.~Scheucher, K.~Swanepoel, P.~Valtr,
\emph{Almost-equidistant sets}, Graphs and Combinatorics \textbf{36}
(2020), 729--754. arXiv:1706.06375.

\bibitem{LR} D.~G.~Larman, C.~A.~Rogers, \emph{The realization of
distances within sets in Euclidean space}, Mathematika \textbf{19} (1972),
1--24.

\bibitem{KMS} A.~Kupavskii, N.~Mustafa, K.~Swanepoel, \emph{Bounding the
size of an almost-equidistant set in Euclidean space}, Combinatorics,
Probability and Computing \textbf{28} (2019), 280--286.

\bibitem{Gyorey} B.~Gy\"orey, \emph{Diszkr\'et metrikus terek
be\'agyaz\'asai}, Master's thesis, E\"otv\"os Lor\'and University,
Budapest, 2004.

\bibitem{BGS} G.~Brinkmann, J.~Goedgebeur, J.-C.~Schlage-Puchta,
\emph{Ramsey numbers $R(K_3,G)$ for graphs of order 10}, Electron.\ J.
Combin.\ \textbf{19} (2012), \#P36. Software \texttt{triangleramsey},
\url{https://caagt.ugent.be/triangleramsey/}.

\bibitem{BBH} S.~Brandt, G.~Brinkmann, T.~Harmuth, \emph{The generation of
maximal triangle-free graphs}, Graphs and Combinatorics \textbf{16}
(2000), 149--157.
\end{thebibliography}

\end{document}
